\section{Giving the Graph a \code{node} Field}
\label{sec:ch13-giving-the-graph-a-node-field}
\index{node@\code{Query.node}!in a federated graph|(}%

Read the composed schema at tag \code{ch12} and you find this:

\begin{graphqlsdl}[fontsize=\footnotesize]
type Product implements Node { ... }
type Customer implements Node { ... }
type Review implements Node { ... }
type Order implements Node { ... }
\end{graphqlsdl}

Four types promising that a client can hold their identifier and come back for
them. Then grep the same document for a field that returns a \code{Node} and
there is not one. The interface is in the schema, four types implement it,
nothing can reach it.

\index{registerNodeInterface@\code{registerNodeInterface}!what it drops}%
That is chapter~\ref{ch:the-first-cut-extracting-catalog}'s bill arriving. All
six services reach the global identification setting through the shared
\code{AddMosaicSubgraph}, and until this chapter it was hard-coded there to
register no node interface at all. That keeps the identifier format, the
interface, the \code{implements Node} clauses and every node resolver, and drops
exactly two root fields. It had to: two subgraphs declaring \code{Query.node} is
a composition error, so a service that turned it back on broke everybody's build
rather than its own. The feature left and the promise stayed, which is the worst
of the available arrangements and is what this section is for.

\subsection*{Two identities, and only one of them travels}

\index{Node interface@\code{Node} interface!versus a federation key}%
Before building anything it is worth being precise about what \code{Node}
promises, because Mosaic has two notions of identity and I had been treating them
as one.

\code{Node} says: a client may hold this identifier and ask the graph for the
thing again. \code{@key} says: another service may hand this identifier to the
router and have the thing found. They look identical from inside one process,
where both mean \enquote{there is a primary key}. Across a network they are
different promises kept by different machinery, and Mosaic kept the first for
four types and the second for two.

\begin{minted}[fontsize=\footnotesize]{text}
                implements Node   @key(fields: "id")
Product              yes               yes
Customer             yes               yes
Review               yes               no
Order                yes               no
\end{minted}

A \code{node} field that can return a \code{Review} needs the router to be able
to fetch one from an identifier, and nothing in the graph could. So the first
part of the work was not about \code{node} at all: \code{Review} and
\code{Order} became entities, with a \code{[Key("id")]} on the record and a
reference resolver behind the DataLoader their node resolvers were already
using. Ordering had never declared a resolvable key before, so its published
schema grew an \code{\_entities} field and a one-member \code{\_Entity} union.

\subsection*{The service that owns the field}

\index{Mosaic.Nodes@\code{Mosaic.Nodes}!the seventh subgraph}%
Somebody has to publish \code{node}, and whoever does it has to be able to return
any type the graph considers addressable, because a resolver can only return a
type its own schema declares. That is real coupling and there is no version of
this feature without it. The choice is where to put it.

Putting it in Catalog would mean the product service declaring \code{Customer},
\code{Review} and \code{Order}, which is exactly the knowledge
chapter~\ref{ch:strangling-the-monolith} spent a chapter removing from one
process. So Mosaic gets a seventh service instead, on 5107, and it owns nothing
else. Its published schema is sixty lines. Below is everything in it that is not
the \code{Node} interface, \code{\_Service}, the federation scalars and the two
directive definitions every subgraph in this book carries:

\begin{graphqlsdl}[fontsize=\footnotesize]
type Query {
  "Fetches an object given its ID."
  node("ID of the object." id: ID!): Node
  "Lookup nodes by a list of IDs."
  nodes("The list of node IDs." ids: [ID!]!): [Node]!
  _service: _Service!
}

type Customer implements Node @key(fields: "id", resolvable: false) {
  id: ID!
}

type Order implements Node @key(fields: "id", resolvable: false) {
  id: ID!
}
\end{graphqlsdl}

\code{Product} and \code{Review} are the same three lines with a different name.

No \code{\_entities}, because every key is \code{resolvable: false} and nothing
here is an entity this service answers for. No database, no connection string and
no \code{depends\_on} in the compose file. It is the smallest service in the
repository and the only one whose file listing is a description of the whole
graph rather than of a domain.

\index{node resolver!handing the key back}%
The resolvers are two lines each and there is a nice surprise in them:

\begin{minted}[fontsize=\footnotesize,breaklines]{csharp}
[NodeResolver]
public static Product ResolveProduct(Guid id) => new() { Id = id };
\end{minted}

Nothing in this project parses a base64 string. HotChocolate's \code{node} field
decodes the identifier itself, reads the type name out of it, finds the node type
of that name and calls its resolver with the internal identifier. Four copies of
\code{ProductKey.TryDecode} exist elsewhere in this repository, and two more of
exactly the same shape for \code{Review} and \code{Order}, because a
\emph{reference} resolver is handed its key raw, and the federation package has
never heard of relay identifiers. The
\code{node} field is the one place in the stack that has, because global object
identification is HotChocolate's own convention rather than federation's. The
two mechanisms meet here and only one of them needs telling what the string is.

\subsection*{What comes back}

Figure~\ref{fig:ch13-node-dispatch} is the shape of it, and the shape is what
makes the feature worth the service.

\begin{figure}[htbp]
  \centering
  % How one identifier reaches four owners. Referenced from chapter 13. Bare
% tikzpicture; the figure environment, caption and label live at the call site.
%
% The point of the picture is the shape of the first hop. Everything the node
% service returns is the type name and the key it was already given, and every
% field the client actually wanted comes out of a second round of fetches that
% could not start until it answered.
\begin{tikzpicture}[
    font=\small,
    box/.style={draw, rounded corners=2pt, minimum width=18mm,
                minimum height=6mm, align=center, font=\scriptsize},
    owner/.style={box, minimum width=16mm},
    par/.style={draw, dashed, rounded corners=2pt, gray},
    flow/.style={-{Stealth[length=2mm]}},
    note/.style={font=\scriptsize, gray, align=center, inner sep=1pt}
  ]

  \node[box] (client) at (0.6,0) {client};
  \node[box] (router) at (4.8,0) {router};
  \node[box, fill=black!8] (nodes) at (9.0,0) {\code{nodes}};

  \draw[flow] (client) -- node[note, above, align=center]
    {\code{node(id:)}} (router);
  \draw[flow] (router) -- node[note, above, align=center]
    {\code{\_\_typename}\\and the key} (nodes);

  \node[note, anchor=north, align=center] at (9.0,-0.55)
    {no database,\\no domain};

  % -- second round -----------------------------------------------------------
  \draw[par] (1.5,-3.0) rectangle (8.1,-2.1);
  \node[owner] (cat) at (2.6,-2.55) {\code{catalog}};
  \node[owner] (pri) at (4.8,-2.55) {\code{pricing}};
  \node[owner] (inv) at (7.0,-2.55) {\code{inventory}};

  \draw[flow] (router) .. controls (4.8,-1.4) and (2.6,-1.4) .. (cat);
  \draw[flow] (router) .. controls (4.8,-1.4) and (4.8,-1.4) .. (pri);
  \draw[flow] (router) .. controls (4.8,-1.4) and (7.0,-1.4) .. (inv);

  \node[note, anchor=west, align=left] at (8.3,-2.55)
    {\code{\_entities}, in\\parallel, keyed on\\what came back};

  \node[note, anchor=east, align=right] at (1.3,-2.55) {title, price,\\quantity};

\end{tikzpicture}

  \caption{One identifier, and the three services that answer for a product.
  The node service replies with the type name and the key it was handed, which
  is all it has; every field the client actually asked for arrives in a second
  round of entity fetches that could not begin until the first hop said what
  kind of thing this was. The service in the middle knows the whole graph's type
  list and none of its data, which is the trade the feature is.}
  \label{fig:ch13-node-dispatch}
\end{figure}

Asking for all four types, with fields the node service does not have:

\begin{minted}[fontsize=\scriptsize,breaklines]{text}
Product:  { "__typename": "Product", "title": "Beech Cutting Board",
            "price": { "amount": 69 }, "availableQuantity": 61 }
Review:   { "__typename": "Review", "rating": 5,
            "author": { "displayName": "Katharina Brandt" } }
Customer: { "__typename": "Customer", "displayName": "Katharina Brandt",
            "email": "k.brandt@example.org" }
Order:    { "__typename": "Order", "placedAt": "2026-04-07T19:20:00Z",
            "total": { "amount": 858 }, "lines": [ { "quantity": 2 } ] }
\end{minted}

Four services answered the first of those and none of them was asked for by
name. The plan says how:

\begin{minted}[fontsize=\scriptsize,breaklines]{text}
Sequence
  Single  nodes    node(id: $a){ __typename ... on Product { __typename id } }
  Parallel
    Entity catalog  dependsOnFetchIds [0]   ... on Product { __typename title }
    Entity pricing  dependsOnFetchIds [0]   ... on Product { price { amount } }
\end{minted}

\code{nodes(ids:)} works the same way and answers a mixed list positionally, so a
client that asked for a product, a review and a customer gets three objects back
in that order.

\subsection*{What it cannot do, and what it must not do}

\index{node@\code{Query.node}!an identifier for a row nobody has}%
Hand it a well-formed identifier for a row that does not exist:

\begin{minted}[fontsize=\footnotesize,breaklines]{text}
Cannot return null for non-nullable field 'Query.node.title'.
\end{minted}

The node service cannot check anything. It decodes a string, finds a type name it
recognises, and returns a stub, because it has nothing to look the identifier up
in. The failure surfaces two hops later at the owner, as a non-null violation on
a field the client happened to select. That is the honest cost of the design and
it is not fixable from inside it: the only service that could answer
\enquote{does this exist} is the one that owns the data, and asking it would make
this service a router.

The failures it does catch, it catches well:

\begin{minted}[fontsize=\footnotesize,breaklines]{text}
There is no node resolver registered for type `OrderLine`.
The node ID string has an invalid format.
\end{minted}

\code{OrderLine} is a type in the graph and is deliberately not a \code{Node}: a
line has no identity outside the order it belongs to. There is no stub for it in
the node service and there never will be, and the error names the type.

\index{shareable@\code{@shareable}!on \code{Query.node}}%
Then there is the arrangement to avoid, and it is the one that composes.
Chapter~\ref{ch:the-first-cut-extracting-catalog} asserted that marking
\code{Query.node} \code{@shareable} in two subgraphs \enquote{would be a lie}
without measuring it. Without the directive, the composer refuses:

\begin{minted}[fontsize=\footnotesize,breaklines]{text}
The Object "Query" defines the same fields in multiple subgraphs without the "@shareable" directive:
 The field "node" is defined in the following subgraphs: "catalog", "nodes".
 However, it is not declared "@shareable" in any of them.
\end{minted}

Add it to both and the graph composes, and the routing table is the finding:

\begin{minted}[fontsize=\footnotesize]{text}
Query.node    catalog, nodes
\end{minted}

Two services advertised as able to answer, a router free to pick either, and one
of them able to resolve one type out of four. \code{@shareable} is a promise
that every service declaring the field gives the same answer, and there is no
error for making that promise about two services that do not. The build goes green and half the identifiers in the graph
start answering null.

\index{key@\code{@key}!makes a field implicitly shareable}%
One more thing about those stubs, because it explains why they are legal at all.
Drop the \code{@key} from the four of them, leaving \code{implements Node} and an
\code{id}, and the composer produces four errors:

\begin{minted}[fontsize=\footnotesize,breaklines]{text}
The Object "Product" defines the same fields in multiple subgraphs without the "@shareable" directive:
 The field "id" is defined and declared "@shareable" in the following subgraphs: "catalog", "pricing", "inventory",
"reviews", "ordering".
 However, it is not declared "@shareable" in the following subgraph: "nodes".
\end{minted}

A key field is implicitly shareable and an ordinary field is not. That is the
only reason five services can all declare \code{Product.id} without saying
anything about it, and it is why the same two-line stub is unremarkable with the
directive and a build failure without it.

\medskip
\code{Node} is an interface, and interfaces are the next thing that stops being
simple at a boundary.

\index{node@\code{Query.node}!in a federated graph|)}%
